\introduction  %% \introduction[modified heading if necessary]
\label{sec:intro}


1.1 Scientific and measurement problem

- Establish the role of OCO-2 XCO2 in flux inversions, regional carbon
  budgets, trend monitoring, and emission studies.
- Explain why sub-ppm systematic errors matter.
- Introduce cloud-side illumination, shadowing, and altered photon paths as
  mechanisms affecting nominally clear soundings near clouds.
- Reframe screening as spatially correlated sampling, not merely reduced
  sample size: persistently cloudy regions lose observations preferentially.

Suggested agency-level sentence:

> Cloud screening is therefore not only a loss of sample size: it
> preferentially removes observations from persistently cloudy regions,
> potentially reinforcing spatial sampling biases in atmospheric CO2
> inversions.

1.2 Previous approaches and unresolved gap

Organize the literature by approach:

1. 3-D radiative-transfer characterization of cloud adjacency;
2. MODIS- or imager-dependent empirical corrections;
3. general retrieval bias correction and machine learning;
4. the operational OCO-2 bias-correction chain.

Define the gap narrowly: no existing approach is shown here to combine
physical interpretation, single-footprint inference without an imager,
probabilistic uncertainty, independent validation, and explicit
plume-preservation tests.

Avoid claiming that the operational B11 correction is not cloud-aware in any
sense. The supported statement is that it does not fully resolve the observed
surface-dependent near-cloud residuals and can over-correct near-cloud land
soundings in this evaluation.

1.3 Questions and contributions

Frame the study around four questions:

1. Do nominally clear OCO-2 spectra contain a measurable signature of nearby
   clouds?
2. Can photon path-length statistics explain the land–ocean difference in
   XCO2 response?
3. Can a per-footprint model reduce independent-reference error without cloud
   distance at inference?
4. Does the correction preserve real CO2 enhancements, and where does it
   fail?

End with five contributions:

- spectrum-derived photon path-length diagnostics;
- a surface-specific probabilistic correction;
- date-blocked, leakage-guarded evaluation;
- independent TCCON, ATom, and shipborne validation;
- smoother, feature-ablation, plume-preservation, and failure-mode tests.

State deployment independence here. Reserve the full three-tier
imager-independence argument for the Discussion.

